\documentclass[11pt]{article}
\usepackage[a4paper,margin=2.4cm]{geometry}
\usepackage{amsmath,amssymb}
\usepackage{graphicx}
\usepackage{booktabs}
\usepackage{siunitx}
\usepackage[hidelinks]{hyperref}
\usepackage{microtype}
\graphicspath{{figs/}}

\newcommand{\Sab}{S_{\mathrm{ab}}}
\newcommand{\Sinc}{S_{\mathrm{inc}}}
\newcommand{\nhat}{\hat{\mathbf{n}}}
\newcommand{\khat}{\hat{\mathbf{k}}}
\newcommand{\erf}{\ensuremath{\operatorname{erf}}}
\newcommand{\ReLU}{\ensuremath{\operatorname{ReLU}}}

\title{Self-shadowing visibility: exactness, a JAX Faddeeva, and a verification suite}
\author{AEGIS internal note}
\date{2026-06-07}

\begin{document}
\maketitle

\begin{abstract}
The self-shadowing design replaces the binary occlusion factor with a single
differentiable gate $V=\tfrac12[1+\erf(c/\sigma)]$ on a baked signed angular
clearance $c$. This note answers three questions about it. First, how much
exactness do we lose against true knife-edge diffraction, and what would it cost
to recover. The exact transition is a Fresnel integral, that is an error
function of a complex argument, so we implement the Faddeeva function in JAX
(Weideman rational approximation) and verify it to a relative error of
\num{1.1e-13} against SciPy with the knife-edge gate accurate to \num{8.9e-15}
and the geometric-boundary value exactly $1/4$ ($-6$\,dB). It costs
$2.9\times$ a real \erf{} per element with the same $O(N)$ scaling
(measured exponent $0.89$ versus $0.90$), so it never changes the computational
scaling. Second, the small-angle step $\sin c\approx c$ is unnecessary and we
use $\sin c$ exactly at no cost; the approximation it replaces was below
\SI{4.3}{\percent} even at the int8 clearance cap. Third, we give the
verification suite (numerical, physical, differentiable, performance) that pins
the regime equivalences and guards against regressions.
\end{abstract}

\section{The gate}

For an incidence direction $\khat$ and outward normal $\nhat$, with
$\mu=\nhat\cdot(-\khat)$ the incidence cosine, the absorbed power density is
\begin{equation}
  \Sab(\mathbf{r}) \;\propto\; g(\mu)\;\cdot\;V\!\left(c(\mathbf{r},\khat)\right),
  \qquad
  V(c)=\tfrac12\!\left[1+\erf\!\left(\frac{c}{\sigma}\right)\right].
  \label{eq:gate}
\end{equation}
$g$ is the existing local activation ($\mu_+$, or the diffraction GELU
$\mu\,\tfrac12[1+\erf(\mu/\sigma_{\mathrm{loc}})]$ with
$\sigma_{\mathrm{loc}}=\sqrt{\lambda/2\pi R_j}$). $V$ is the new distal factor.
$c$ is the signed geodesic angular clearance from $\khat$ to the nearest distal
occlusion boundary on the direction sphere ($c>0$ visible). The penumbra width
follows from knife-edge geometry: a ray clears the binding edge (distance $d_2$)
by $h=d_2\sin c$, and with the Fresnel parameter
$\nu=h\sqrt{(2/\lambda)(1/d_1+1/d_2)}$,
\begin{equation}
  \sigma \;=\; K\sqrt{\frac{\lambda}{d_{\mathrm{occ}}}},\qquad
  K=\frac{1}{\sqrt{2\pi}},
  \label{eq:sigma}
\end{equation}
the same $\sqrt{\lambda/\text{distance}}$ form as the local GELU, with
$d_{\mathrm{occ}}=R_j$ locally and $d_2$ distally (the $d_1$ term recovers the
near-field width when the source is at finite range). Figure~\ref{fig:penumbra}
shows the width across FR1 and mmWave.

\begin{figure}[t]
  \centering
  \includegraphics[width=0.62\linewidth]{penumbra_width}
  \caption{Distal gate versus angular clearance at three frequencies,
  $d_{\mathrm{occ}}=\SI{0.2}{m}$. The penumbra narrows as $\sqrt{\lambda}$, from
  $\sigma\approx\ang{33}$ at \SI{0.7}{GHz} to $\ang{5}$ at \SI{28}{GHz}.}
  \label{fig:penumbra}
\end{figure}

\section{How exact is the \erf{} gate, and what would exactness cost}

\subsection{The exact transition is a complex error function}

True knife-edge diffraction gives the field relative to free space as the
Fresnel integral
\begin{equation}
  F(\nu)=\frac{1+i}{2}\int_\nu^\infty e^{-i\pi t^2/2}\,dt
        =\frac{1+i}{2}\Big[\big(\tfrac12-C(\nu)\big)-i\big(\tfrac12-S(\nu)\big)\Big],
\end{equation}
and the Fresnel integrals are an error function of a complex argument,
\begin{equation}
  C(\nu)+iS(\nu)=\frac{1+i}{2}\,\erf\!\left(\frac{\sqrt\pi}{2}(1-i)\,\nu\right).
  \label{eq:fresnel-erf}
\end{equation}
So ``exact'' is not a different kind of computation, it is $\erf$ of a complex
argument, equivalently the Faddeeva function $w(z)=e^{-z^2}\operatorname{erfc}(-iz)$
through $\erf(z)=1-e^{-z^2}w(iz)$. SciPy provides this; JAX ships only a real
\erf, so the exact path under the JAX backend needs a custom Faddeeva.

\subsection{A JAX Faddeeva, verified}

We implement $w(z)$ with Weideman's rational approximation (48 terms): the
coefficients are computed once in NumPy, the evaluation is pure
\texttt{jax.numpy}, vectorised, jittable and differentiable, with an
upper-half-plane evaluation and the reflection $w(-z)=2e^{-z^2}-w(z)$ for
$\operatorname{Im}z<0$. Verified against \texttt{scipy.special.wofz} on a
$240\times240$ grid over $[-6,6]^2$ (Fig.~\ref{fig:faddeeva}), the maximum
relative error is \num{1.1e-13}. The derived knife-edge gate $|F(\nu)|^2$ matches
\texttt{scipy.special.fresnel} to \num{8.9e-15} absolute, with the
geometric-boundary value $|F(0)|^2$ equal to $1/4$ to machine precision (the
$-6$\,dB point).

\begin{figure}[t]
  \centering
  \includegraphics[width=0.6\linewidth]{faddeeva_error}
  \caption{Relative error of the JAX Faddeeva against SciPy over the complex
  plane. Below \num{1e-13} everywhere; the structure is round-off, not method
  error.}
  \label{fig:faddeeva}
\end{figure}

\subsection{Cost: a constant factor, not a scaling change}

Table~\ref{tab:timing} and Fig.~\ref{fig:timing} settle the cost question. The
JAX Faddeeva costs \SI{29.7}{ns} per element versus \SI{10.3}{ns} for a real
\erf, a factor of $2.9$, and is $5.5\times$ faster than SciPy's \texttt{wofz}.
All three scale as $O(N)$ (fitted log-log exponents $0.89$, $0.90$, $1.00$): the
exact gate changes the constant, never the scaling. This matches the working
assumption that a fixed multiple is acceptable as long as the asymptotics are
untouched.

\begin{table}[t]
  \centering
  \caption{Per-element cost and measured $O(N^p)$ scaling exponent
  (CPU, float64, $N$ up to \num{1e7}).}
  \label{tab:timing}
  \begin{tabular}{lcc}
    \toprule
    routine & ns/element & scaling exponent $p$ \\
    \midrule
    real \erf{} (JAX, jit)        & 10.3 & 0.90 \\
    Faddeeva (JAX, jit)           & 29.7 & 0.89 \\
    \texttt{wofz} (SciPy)         & 164  & 1.00 \\
    \bottomrule
  \end{tabular}
\end{table}

\begin{figure}[t]
  \centering
  \includegraphics[width=0.6\linewidth]{timing}
  \caption{Wall time versus array length. Parallel log-log lines confirm a
  common $O(N)$ scaling; the JAX Faddeeva sits a constant $\sim\!3\times$ above a
  real \erf{} and below SciPy.}
  \label{fig:timing}
\end{figure}

\subsection{What the erf surrogate actually loses}

Figure~\ref{fig:family} compares the three gates in the dimensionless clearance
$s=c/\sigma$. The default unsquared $\tfrac12[1+\erf s]$ sits at $1/2$ at the
boundary ($-3$\,dB), the squared form $[\tfrac12(1+\erf s)]^2$ matches the exact
$1/4$ ($-6$\,dB), and the full $|F|^2$ adds the lit-side Cornu ripples. The
ripples are precisely what incoherent summation and \SI{4}{cm^2} averaging
remove, so they do not justify the Faddeeva in production. The one real loss is
the shadow tail (Fig.~\ref{fig:tail}): the exact gate decays as $\nu^{-2}$ while
the \erf{} forms decay as $e^{-s^2}$, so deep-shadow leakage is under-predicted,
by a factor of \num{2.1e5} at $s=4$. This is the same defect as the
single-edge approximation: in deep shadow multiple edges leak additively, and
we keep only the binding edge.

\begin{figure}[t]
  \centering
  \includegraphics[width=0.6\linewidth]{gate_family}
  \caption{Gate family in $s=c/\sigma$. Default \erf{} ($-3$\,dB at $s=0$),
  squared \erf{} ($-6$\,dB, knife-edge-consistent), and the exact $|F|^2$ with
  lit-side ripples.}
  \label{fig:family}
\end{figure}

\begin{figure}[t]
  \centering
  \includegraphics[width=0.6\linewidth]{shadow_tail}
  \caption{Deep-shadow leakage. The exact $\nu^{-2}$ tail is many orders above
  the Gaussian \erf{} decay; the monotone surrogate under-predicts dose reaching
  fully shadowed regions.}
  \label{fig:tail}
\end{figure}

\paragraph{Recommendation.} Keep the real-\erf{} gate as the default. The
$-6$\,dB boundary is free via the squared form (\texttt{gate="knife\_edge"}); the
shadow tail and additive multi-edge are a single scoped refinement for
deep-shadow dosimetry (a $\nu^{-2}$ tail graft, or the $k$ nearest edges). The
Faddeeva is now available and verified should a fully exact mode ever be wanted,
at $2.9\times$ and no scaling change.

\section{The $\sin c\approx c$ question: use it exactly}

The clearance-to-clearance map $h=d_2\sin c$ is exact geometry; only the
linearisation $\sin c\approx c$ was an approximation, and it is unnecessary. We
use $\sin c$ directly (a one-token change, no cost). For the record,
Fig.~\ref{fig:sinc} shows the error the linearisation would have carried: it
enters only the penumbra \emph{width} (second order), never the boundary
\emph{location} ($\sin 0=0$ exactly), and stays below \SI{4.3}{\percent} even at
the int8 clearance cap of \SI{0.5}{rad}, well inside the typical $\sigma$ at FR1
and mmWave. So the question is moot: exact is free, and even the approximation
was negligible.

\begin{figure}[t]
  \centering
  \includegraphics[width=0.6\linewidth]{sinc_error}
  \caption{Relative error of $\sin c\approx c$. Marked: the int8 clearance cap
  and the penumbra widths $\sigma$ at \SI{2.45}{GHz} and \SI{28}{GHz}
  ($d_{\mathrm{occ}}=\SI{0.2}{m}$). We avoid it entirely by keeping $\sin c$.}
  \label{fig:sinc}
\end{figure}

\section{Verification suite: make it work and keep it working}

Table~\ref{tab:tests} is the test catalogue. The numerical group pins the proved
regime equivalences (the gate reduces exactly to today's kernels when shadowing
is off, convex, or in the geometric limit). The physical group anchors the new
behaviour to independent references. The differentiable group protects the JAX
gradient path the near-field design loop depends on. The performance group is a
regression guard so the $O(N)$ query and the bounded bake do not silently
degrade.

\begin{table}[t]
  \centering
  \caption{Verification and regression suite for the visibility gate.}
  \label{tab:tests}
  \small
  \begin{tabular}{p{0.30\linewidth}p{0.40\linewidth}p{0.22\linewidth}}
    \toprule
    test & guards & reference \\
    \midrule
    \multicolumn{3}{l}{\emph{numerical (regime equivalence)}}\\
    octahedral round-trip & encode/decode identity, seam continuity & analytic \\
    clearance vs great-circle & distance transform correctness & brute force \\
    convex $\Rightarrow V\equiv1$ & Mie canary, P2 & monograph $\eta\equiv1$ \\
    $\sigma\to0\Rightarrow$ binary & geometric limit, P4 & \texttt{compute\_visibility} oracle \\
    local limit bit-for-bit & recovers ReLU/GELU, P3, F1 & \texttt{physical\_gelu} \\
    opt-out byte-identity & \texttt{self\_shadow=False}, P1 & pre-feature baseline \\
    multi-path covariance & per-path fold vs old mean, P5 & analytic counterexample \\
    \midrule
    \multicolumn{3}{l}{\emph{physical}}\\
    $A_{\mathrm{ab}}=\int\eta$ & exposure-area scaling & monograph $0.80/0.87$ \\
    penumbra $\propto\sqrt\lambda$ & width law, Eq.~\eqref{eq:sigma} & Fig.~\ref{fig:penumbra} \\
    asymmetric-occluder sign & direction convention, F6 & one-sided slab \\
    knife-edge gate accuracy & exact-mode correctness & \texttt{scipy.fresnel} (\num{9e-15}) \\
    \midrule
    \multicolumn{3}{l}{\emph{differentiable}}\\
    Faddeeva accuracy & JAX complex \erf{} & \texttt{scipy.wofz} (\num{1.1e-13}) \\
    near-field $d_1\le0$ guard & no NaN value or gradient, F3 & double-where \\
    autodiff vs finite-diff & pose gradient, $<\SI{2}{\percent}$ & finite difference \\
    \midrule
    \multicolumn{3}{l}{\emph{performance}}\\
    query scaling & $O(N)$ lookup, no ray cast at query & exponent $\approx1$ \\
    bake budget & tens of seconds per phantom, cached & wall-clock bound \\
    \bottomrule
  \end{tabular}
\end{table}

\section{Conclusions}

The \erf{} gate is the right default: it reduces exactly to the current kernels
in the off, convex, and geometric limits, and the only physics it drops (lit
ripples) is averaged away. Exactness is now within reach and verified, the JAX
Faddeeva matching SciPy to \num{1.1e-13} at $2.9\times$ a real \erf{} and the
same $O(N)$ scaling, so a fully exact mode is a constant-factor option rather
than a scaling risk. The two refinements worth keeping on the roadmap are the
$-6$\,dB squared boundary (free) and a deep-shadow tail plus additive multi-edge
(for deep-shadow dosimetry only). The $\sin c$ map is used exactly. The suite in
Table~\ref{tab:tests} makes the equivalences and the differentiability testable
and keeps them from regressing.

\end{document}
